\section{El Problema}
Los hermanos siguieron creciendo. Mateo también aprendió sobre programación dinámica, y cada uno aplicaba la lógica sabiendo que el otro también lo hacía. El juego de las monedas se tornó aburrido en cuanto notaron que siempre ganaba quien empezara, o según la suerte. Los años pasaron, llegó la adolescencia y empezaron a tener gustos diferentes. En general, jugaban a juegos individuales. En particular, Sophia estaba muy enganchada con el siguiente juego de tipo \textit{puzzle}: El juego consiste de una grilla de $n \times m$ celdas. Cada celda puede estar vacía, estar ocupada por una piedra verde, o estar ocupada por una piedra celeste. El objetivo original del juego es remover algunas de las piedras de tal forma que las piedras que quedan cumplan las siguientes condiciones:

\begin{enumerate}
    \item Cada fila tiene al menos una piedra.
    \item Ninguna columna tiene piedras de ambos colores.
\end{enumerate}

A Sophia le interesa tanto esta versión, como también aquella que \textbf{relaja} la primera restricción, obteniendo una solución que maximice la cantidad de fila que cumplen la restricción (manteniendo siempre la segunda restricción). Es decir, minimizar la cantidad de filas que quedan vacías.


Para transformar esta situación en un problema algorítmico concreto, el desafío es el siguiente:
Hacer un análisis del problema, y cumplir varios puntos (después explicados) del problema:
Dado un tablero de $n \times m$ celdas, donde cada celda está vacía u ocupada con una piedra verde o celeste (solo puede estar ocupada por 1 como mucho), ¿es posible remover piedras de tal forma que cada fila tenga al menos una piedra y cada columna no tenga piedras de ambos colores?

Los puntos a resolver son:
\begin{itemize}
    \item Demostrar que el Problema del juego de las Piedras se encuentra en NP, y es un problema \textit{NP-Completo}.
    \item Escribir para obtener la solución óptima al problema en su versión de optimización: un algoritmo de backtracking y un modelo de programación lineal
    \item Analizar el famoso algoritmo de aproximación de ``El Gringo Hinz'' (\ref{sec:hinz})
    \item Implementar un algoritmo de aproximación para el problema.
\end{itemize}
\subsection{Introducción}

Tranquilamente se pueden observar varios puntos a realizar. Lo ideal, como primera instancia, hay que tener en cuenta que existen tableros \textbf{sin} solución, los cuales no se pueden cumplir las dos condiciones iniciales del problema. A eso sumémosle que si se demuestra que este problema es NP-Completo, estamos frente a un problema \textit{muy difícil} de resolver.
Nos podemos preguntar:
\begin{itemize}
    \item ¿Es realmente posible encontrar una solución óptima? Si tanto se habla de que es muy difícil, ¿\textit{cuánto tardaría}?
    \item ¿Cuáles son las diferencias entre un algoritmo de Backtracking y Programación Lineal? ¿Distan mucho de los algoritmos de Aproximación? ¿Todos me aseguran optimalidad?
    \item ¿Cuánto me beneficia el Problema del Juego de las Piedras en su \textit{versión de Relajación}?
\end{itemize}

Una vez propuesto eso, sigamos pensando para ahora empezar con el planteo.

\subsection{Planteamiento y Modelado del problema}

Para cualquier tablero que se nos presente, nuestro objetivo va a ser ganar siempre que podamos.
En la Figura \ref{fig:con_solucion} se observa una instancia válida para el problema de optimización

\begin{figure}[H]
    \centering
    % tablero 3x3 con solución en versión de decisión
\begin{tikzpicture}[
    scale=0.85,
    transform shape,
    celda/.style={draw=gray!60, fill=gray!10, line width=0.6pt},
    verde/.style={circle, fill=green!60!black, draw=green!40!black, minimum size=0.68cm, inner sep=0pt},
    celeste/.style={circle, fill=cyan!75!blue!80, draw=cyan!50!black, minimum size=0.68cm, inner sep=0pt},
    tachado/.style={circle, draw=red!80!black, line width=1.3pt, densely dashed, minimum size=0.72cm, inner sep=0pt}
]

% macro local para dibujar la cruz roja sobre una piedra
\newcommand{\cruz}[2]{
    \node[tachado] at (#1, #2) {};
    \draw[red!80!black, line width=1.4pt] (#1-0.22, #2-0.22) -- (#1+0.22, #2+0.22);
    \draw[red!80!black, line width=1.4pt] (#1-0.22, #2+0.22) -- (#1+0.22, #2-0.22);
}

%  TABLERO INICIAL
\begin{scope}[shift={(0,0)}]
    \draw[celda] (0,0) rectangle (3,3);
    \draw[celda] (0,0) grid (3,3);
    
    % Fila 1 (arriba: y=2.5)
    \node[verde] at (0.5, 2.5) {};
    \node[celeste] at (1.5, 2.5) {};
    % Fila 2 (medio: y=1.5)
    \node[verde] at (0.5, 1.5) {};
    \node[celeste] at (2.5, 1.5) {};
    % Fila 3 (abajo: y=0.5)
    \node[celeste] at (1.5, 0.5) {};
    \node[verde] at (2.5, 0.5) {};
    
    \node[font=\footnotesize\sffamily] at (1.5, -0.45) {Tablero original};
\end{scope}

%  PIEDRAS A REMOVER 
\begin{scope}[shift={(4.2,0)}]
    \draw[celda] (0,0) rectangle (3,3);
    \draw[celda] (0,0) grid (3,3);
    
    % Fila 1
    \node[verde] at (0.5, 2.5) {};
    \node[celeste] at (1.5, 2.5) {}; \cruz{1.5}{2.5}
    % Fila 2
    \node[verde] at (0.5, 1.5) {}; \cruz{0.5}{1.5}
    \node[celeste] at (2.5, 1.5) {};
    % Fila 3
    \node[celeste] at (1.5, 0.5) {};
    \node[verde] at (2.5, 0.5) {}; \cruz{2.5}{0.5}
    
    \node[font=\footnotesize\sffamily] at (1.5, -0.45) {Piedras removidas};
\end{scope}

% RESULTADO VÁLIDO 
\begin{scope}[shift={(8.4,0)}]
    \draw[celda] (0,0) rectangle (3,3);
    \draw[celda] (0,0) grid (3,3);
    
    % Fila 1 (Col 1: Verde)
    \node[verde] at (0.5, 2.5) {};
    % Fila 2 (Col 3: Celeste)
    \node[celeste] at (2.5, 1.5) {};
    % Fila 3 (Col 2: Celeste)
    \node[celeste] at (1.5, 0.5) {};
    
    \node[font=\footnotesize\sffamily] at (1.5, -0.45) {Solución válida};
\end{scope}

\end{tikzpicture}

    \renewcommand{\thefigure}{1}
    \caption{Ejemplo de tablero $3 \times 3$ con solución en la versión de decision.}
    \label{fig:con_solucion}
\end{figure}

Por otro lado, la Figura \ref{fig:sin_solucion} muestra una configuración donde resulta imposible cubrir todas las filas sin generar conflictos de color en las columnas.

\begin{figure}[H]
    \centering
    %La fila 2 obliga a que la columna 1 sea celeste, y la fila 3 obliga a que la columna 2 sea verde. Con esa configuración, en la fila 1 deben eliminarse ambas piedras y queda vacía.

\begin{tikzpicture}[
    scale=0.85,
    transform shape,
    celda/.style={draw=gray!60, fill=gray!10, line width=0.6pt},
    verde/.style={circle, fill=green!60!black, draw=green!40!black, minimum size=0.7cm, inner sep=0pt},
    celeste/.style={circle, fill=cyan!75!blue!80, draw=cyan!50!black, minimum size=0.7cm, inner sep=0pt}
]

\draw[celda] (0,0) rectangle (3,3);
\draw[celda] (0,0) grid (3,3);

% Fila 1 (arriba: y=2.5) -> Verde en col 1, Celeste en col 2
\node[verde] at (0.5, 2.5) {};
\node[celeste] at (1.5, 2.5) {};

% Fila 2 (medio: y=1.5) -> Solo Celeste en col 1 (fuerza col 1 a ser celeste)
\node[celeste] at (0.5, 1.5) {};

% Fila 3 (abajo: y=0.5) -> Solo Verde en col 2 (fuerza col 2 a ser verde)
\node[verde] at (1.5, 0.5) {};

\end{tikzpicture}

    \renewcommand{\thefigure}{1.1}
    \caption{Ejemplo de tablero $3 \times 3$ sin solución para el problema.}
    \label{fig:sin_solucion}
\end{figure}

Finalmente, al relajar la primera restricción (Figura~\ref{fig:relajacion}), una instancia que no admite solución exacta logra maximizar la cantidad de filas cubiertas dejando únicamente una fila vacía.

\begin{figure}[H]
    \centering
    % sin solucion en decision -> optimo en relajacion
\begin{tikzpicture}[
    scale=0.85,
    transform shape,
    celda/.style={draw=gray!60, fill=gray!10, line width=0.6pt},
    verde/.style={circle, fill=green!60!black, draw=green!40!black, minimum size=0.65cm, inner sep=0pt},
    celeste/.style={circle, fill=cyan!75!blue!80, draw=cyan!50!black, minimum size=0.65cm, inner sep=0pt},
    tachado/.style={circle, draw=red!80!black, line width=1.3pt, densely dashed, minimum size=0.7cm, inner sep=0pt}
]

\newcommand{\cruz}[2]{
    \node[tachado] at (#1, #2) {};
    \draw[red!80!black, line width=1.4pt] (#1-0.21, #2-0.21) -- (#1+0.21, #2+0.21);
    \draw[red!80!black, line width=1.4pt] (#1-0.21, #2+0.21) -- (#1+0.21, #2-0.21);
}

% INSTANCIA INFACTIBLE EN DECISIÓN 
\begin{scope}[shift={(0,0)}]
    \draw[celda] (0,0) rectangle (3,4);
    \draw[celda] (0,0) grid (3,4);
    
    % Fila 1 (y=3.5): Verde en col 1, Celeste en col 2
    \node[verde] at (0.5, 3.5) {};
    \node[celeste] at (1.5, 3.5) {};
    % Fila 2 (y=2.5): Solo Celeste en col 1
    \node[celeste] at (0.5, 2.5) {};
    % Fila 3 (y=1.5): Solo Verde en col 2
    \node[verde] at (1.5, 1.5) {};
    % Fila 4 (y=0.5): Verde en col 3
    \node[verde] at (2.5, 0.5) {};
    
    \node[font=\footnotesize\sffamily] at (1.5, -0.45) {Tablero original};
\end{scope}

%  Seleccion de piedras removidas ÓPTIMA
\begin{scope}[shift={(4.2,0)}]
    \draw[celda] (0,0) rectangle (3,4);
    \draw[celda] (0,0) grid (3,4);
    
    % Fila 1 (queda vacía para salvar filas 2 y 3)
    \node[verde] at (0.5, 3.5) {}; \cruz{0.5}{3.5}
    \node[celeste] at (1.5, 3.5) {}; \cruz{1.5}{3.5}
    % Fila 2
    \node[celeste] at (0.5, 2.5) {};
    % Fila 3
    \node[verde] at (1.5, 1.5) {};
    % Fila 4
    \node[verde] at (2.5, 0.5) {};
    
    \node[font=\footnotesize\sffamily] at (1.5, -0.45) {Removidas óptimas};
\end{scope}

% SOLUCIÓN RELAJADA (3/4 FILAS
\begin{scope}[shift={(8.4,0)}]
    \draw[celda] (0,0) rectangle (3,4);
    \draw[celda] (0,0) grid (3,4);
    
    % Fila 1 vacía (y=3.5) - resaltada suavemente en rojo
    \draw[fill=red!8, draw=red!40, dashed] (0,3) rectangle (3,4);
    
    % Fila 2
    \node[celeste] at (0.5, 2.5) {};
    % Fila 3
    \node[verde] at (1.5, 1.5) {};
    % Fila 4
    \node[verde] at (2.5, 0.5) {};
    
    \node[font=\footnotesize\sffamily] at (1.5, -0.45) {Óptimo ($3$ de $4$ filas)};
\end{scope}

\end{tikzpicture}


    \renewcommand{\thefigure}{1.2}
    \caption{Resolución óptima de una instancia del problema en la versión relajada.}
    \label{fig:relajacion}
\end{figure}
Habiendo encontrado un contraejemplo para que si o si tenemos que dejar al menos una fila vacía, surge la pregunta clave para modelar el problema:
¿cómo tomamos las decisiones sin perdernos en las $n \times m$ celdas?

Si analizamos la segunda regla (\textit{ninguna columna puede tener piedras de ambos colores}), notamos que: una vez que decidimos que en una columna $j$ no va a haber piedras celestes, nunca nos va a convenir eliminar ninguna piedra verde de esa misma columna, porque conservar todas las piedras verdes disponibles nunca rompe la restriccion vertical y al mismo tiempo esta maximizando nuestras chances de cubrir filas horizontales.
Por lo tanto, el problema de elegir qué piedras sacar es matemáticamente igual a agarrar exactamente \textbf{$m$ decisiones binarias} (una por cada columna $j \in \{0, \dots, m-1\}$): o la columna $j$ sigue solo con sus piedras verdes (\texttt{V}), o deja solo sus piedras celestes (\texttt{C}).

Por último, es muy importante diferenciar entre las dos variantes del problema:
\begin{itemize}
    \item \textbf{Versión de Decisión:} Dado un tablero de $n \times m$ celdas, ¿es posible remover piedras de forma tal que ninguna columna tenga piedras de ambos colores y todas las filas tengan \textbf{al menos} una piedra ($0$ filas vacías)? (¿es posible dejar a lo sumo $k$ filas vacías?).
    \item \textbf{Versión de Optimización (Relajada):} Dado un tablero de $n \times m$ celdas, encontrar una asignación válida de colores por columna que \textbf{minimice la cantidad de filas que quedan vacías} (o equivalentemente, que maximice la cantidad de filas cubiertas).
\end{itemize}

\subsection{Consecuencias de la optimalidad}

Si se logra pensar y crear estrategias para las cuales siempre que se pueda ganar lo logremos de forma eficaz, eso significa que:
\begin{itemize}
    \item Las indicaciones propuestas fueron óptimas.
    \item Si pasamos las instrucciones a un algoritmo, ese mismo es óptimo.
    \item Si la estrategia no devuelve siempre el óptimo, sino una solucion \textit{cercana} a la óptima, la solución se considera una \textit{aproximada}.
    \item Si logramos resolverlo polinomialmente, el algoritmo se considera rápido y muy eficaz si nos devuelve siempre la solución óptima. En cambio, si logramos resolverlo exponencialmente habría que considerar \textit{sacrificar} cierta optimalidad para que, con tableros de gran magnitud, no tardemos semanas en resolverlo.
\end{itemize}

Las pruebas a realizar de todo esto y mucho más se va a realizar en la lectura que sigue. Se recomienda leer detenidamente cada paso, y si no se entiende, releerlo. 
A eso sumarle el \textbf{pensamiento} de lograr la optimalidad, hasta acá se dijo sutilmente como es éste y su gran significado en las Ciencias de la Computación. A continuación, se explica exáctamente eso: en dónde se encasillan los \textbf{algoritmos polinomiales}, si es posible \textbf{verificar} una posible solución del problema en tiempo polinomial, y la relación de todo esto con los algoritmos no polinomiales. 

¿Es posible resolver el Problema del Juego de las Piedras, si yo \textbf{sé} resolver otro problema \textit{muy parecido} a éste?

